% !TEX program = xelatex
% Оформление: sadekovlar/misis-article-latex, см. LICENSE.
\documentclass{misisarticle}
\usepackage{needspace}

\grnti{}
\asjc{}
\oecd{}
\title{Гиперсетевая генерация низкоранговых адаптеров для эффективной доменной адаптации VLA-моделей}
\misisauthor{А.~А.~Хабибуллин}{кафедра инженерной кибернетики НИТУ «МИСИС»}{Москва, Россия}{}

\setabstract{Рассматривается задача параметрически- и данно-эффективной доменной адаптации моделей vision-language-action (VLA). На основе метода Hyper-LoRA описана схема генерации низкоранговых адаптеров для компактной модели SmolVLA. Гиперсеть порождает поправки к MLP-слоям языковой части замороженного визуально-языкового модуля, а эксперт действий обучается совместно с гиперсетью. Приведены математические описания низкоранговой адаптации, гиперсетевой генерации весов и согласования потоков как обучающего сигнала. Сформулирована постановка обучения на демонстрациях LIBERO и проверки обобщения на LIBERO-Pro. Работа имеет аналитико-теоретический характер; экспериментальная проверка относится к следующему этапу исследования.}
\setkeywords{VLA, доменная адаптация, LoRA, гиперсети, SmolVLA, поведенческое клонирование, flow matching, LIBERO, LIBERO-Pro}
\hypersetup{pdftitle={Гиперсетевая генерация низкоранговых адаптеров для эффективной доменной адаптации VLA-моделей},pdfauthor={А. А. Хабибуллин}}

\begin{document}
\maketitle

\section{Введение}

Перенос идей фундаментальных моделей из обработки естественного языка и компьютерного зрения в робототехнику породил класс робототехнических фундаментальных моделей (RFM) и, в частности, VLA-моделей, которые по изображениям с камер и инструкции на естественном языке напрямую предсказывают низкоуровневые действия манипулятора~\cite{firoozi2023}. Такие модели обучаются поведенческим клонированием на демонстрациях; линия RT-1~\cite{rt1}, RT-2~\cite{rt2}, Open X-Embodiment~\cite{oxe}, Octo~\cite{octo} и OpenVLA~\cite{openvla} последовательно наращивала масштаб данных и общность, а компактная SmolVLA~\cite{smolvla} показала работоспособность подобной архитектуры при существенно меньших вычислительных затратах.

При переносе предобученной VLA в новый домен -- иную обстановку, набор задач или внешний вид объектов -- ее качество может падать из-за доменного разрыва, что требует доменной адаптации. Практически значима адаптация, одновременно параметрически-эффективная (мало обучаемых весов) и данно-эффективная (мало демонстраций). Остроту проблеме придает ограниченная обобщающая способность существующих политик: контролируемые возмущения бенчмарка LIBERO-Pro~\cite{liberopro} выявили существенное падение успешности моделей, показывающих высокие результаты на стандартном LIBERO~\cite{libero}. Это указывает на риск запоминания обучающих конфигураций и делает проверку обобщения ключевым требованием к методам адаптации.

Целью настоящей работы является теоретическая проработка метода гиперсетевой генерации низкоранговых адаптеров (Hyper-LoRA) для параметрически- и данно-эффективной доменной адаптации VLA-моделей, а также предварительная постановка задачи экспериментального исследования. На данном этапе исследование ограничено анализом предметной области, систематизацией методов и обоснованием архитектурных решений; практическая реализация и эмпирическая проверка отнесены к следующему этапу.

\section{Анализ работ}

Опорным для настоящей работы выступает метод Hyper-LoRA~\cite{hyperlora}, в котором гиперсеть, обусловленная контекстом задачи, порождает специфичные для задачи параметры низкоранговой адаптации для робототехнической фундаментальной модели. Контекст задается инструкцией на естественном языке и кодируется энкодером в эмбеддинг фиксированной размерности, после чего выходные головы генератора порождают матрицы адаптера. Обучаемыми остаются параметры гиперсети, а не непосредственно порождаемые матрицы.

Эффективность метода обеспечивают два приема. Первый -- подача индекса слоя RFM как дополнительного входа гиперсети: единый набор выходных голов переиспользуется для разных слоев, а различие между слоями кодируется их индексом. Поскольку основная масса параметров гиперсети сосредоточена в выходных головах, такое переиспользование сокращает ее размер. Второй прием -- dropout контекстного эмбеддинга перед выходными головами, направленный на снижение переобучения генератора на конкретные обучающие задачи~\cite{hyperlora}.

Постановка в опорной работе строится на базе Octo с целевым доменом SIMPLER~\cite{simpler}: гиперсеть обучается на подмножестве задач, адаптеры добавляются к MLP-слоям, а диффузионная голова дополнительно дообучается. Обобщение проверяется без дополнительной настройки на оставшихся задачах домена~\cite{hyperlora}.

\section{Постановка задачи}

Под доменом понимается среда взаимодействия робота, заданная распределением визуальных сцен, набором манипулируемых объектов и пространством достижимых конфигураций манипулятора, вместе с конечным множеством задач, которые в этой среде требуется решать. Отдельная задача специфицируется инструкцией на естественном языке, описывающей желаемую цель манипуляции. Качество политики оценивается долей успешно завершенных эпизодов.

Поведение эксперта формализуется траекториями:
\begin{equation}
  \tau=(o_0,a_0,o_1,a_1,\ldots,o_H,a_H),
\end{equation}
где $o_t$ -- наблюдение в момент $t$, $a_t$ -- действие манипулятора, $H$ -- горизонт эпизода. Инструкция $c$ задается для всей траектории. Каждой задаче сопоставлен набор демонстраций, собранных под управлением эксперта.

Множество задач обозначим $\mathcal{T}$. Демонстрации доступны на обучающем подмножестве $\mathcal{T}_{\mathrm{train}}$, тогда как практический интерес представляет более широкое целевое множество условий $\mathcal{T}_{\mathrm{target}}$. Требуется по демонстрациям из $\mathcal{T}_{\mathrm{train}}$ построить параметрически- и данно-эффективную настройку предобученной VLA и оценить ее перенос на целевые условия, не представленные в обучении.

Пусть $\theta_0$ -- замороженные параметры VLM, $\phi$ -- параметры гиперсети, $\psi$ -- параметры эксперта действий. Требуется обучить $\phi$ и $\psi$, сохранив $\theta_0$ неизменными. Гиперсеть $H_\phi(c,l)$ порождает LoRA-поправку для слоя $l$; в расширенном варианте контекст $c$ дополнительно включает визуальные признаки текущей сцены.

В качестве обучающих данных планируется использовать LeRobot-конверсию LIBERO \texttt{physical-intelligence/libero}~\cite{libero,liberodata}. Она содержит 1693 эпизода на 40 задачах из наборов LIBERO-Spatial, LIBERO-Object, LIBERO-Goal и LIBERO-10. Наблюдения включают изображения внешней и кистевой камер, состояние манипулятора и языковую инструкцию; действия задаются семимерным вектором управления перемещением, вращением и захватом.

Обобщение предполагается оценивать без дообучения на возмущенных условиях LIBERO-Pro~\cite{liberopro}, отдельно для изменений объектов, начальных положений, формулировок инструкций и целевых задач. Качеством управления служит доля успешно завершенных эпизодов.

\section{Описание методов}

Настоящая постановка наследует базовый механизм Hyper-LoRA, но переносит его на SmolVLA. Визуально-языковой модуль замораживается, низкоранговые адаптеры порождаются поверх MLP-слоев его языковой части, а эксперт действий настраивается совместно с гиперсетью. Веса генерируются для каждого примера батча в процессе прямого прохода. Обобщение предполагается проверять на LIBERO-Pro без дообучения на его возмущенных условиях.

Поведенческое клонирование формулирует обучение политики как задачу обучения с учителем: по наблюдению $o$ и инструкции $c$ предсказывается действие $a$, а целевыми метками служат действия эксперта из демонстраций. В общем виде минимизируется ожидаемое расхождение
\begin{equation}
  \mathcal{L}_{\mathrm{BC}}=\mathbb{E}_{(o,c,a)\sim\mathcal{D}}
  \bigl[d(\pi(o,c),a)\bigr].
\end{equation}
Здесь $\mathcal{D}$ -- распределение демонстрационных примеров, $\pi$ -- обучаемая политика, $d$ -- мера расхождения в пространстве действий. Такое обучение не требует взаимодействия со средой на этапе оптимизации.

Простейшая регрессия одного действия с квадратичной ошибкой плохо согласуется с многомодальностью демонстраций: цель может быть достижима несколькими допустимыми способами, а усреднение разных мод приводит к нежелательному «среднему» действию. Поэтому в современных VLA моделируется условное распределение действий, из которого порождаются выборки~\cite{octo,smolvla}.

В качестве генеративного механизма в SmolVLA применяется согласование потоков~\cite{smolvla,flowmatching}. Вводится семейство промежуточных состояний между шумовой выборкой $x_0$ и экспертным блоком действий $a$, параметризованное временем $s\in[0,1]$. Для прямолинейного пути
\begin{equation}
  x_s=(1-s)x_0+sa,\qquad u_s=a-x_0.
\end{equation}
Обучаемое векторное поле $v$ приближает $u_s$ при условии наблюдения и инструкции:
\begin{equation}
\begin{split}
  \mathcal{L}_{\mathrm{FM}}=\mathbb{E}_{o,c,a,x_0,s}\Bigl[
  \bigl\|v_{\theta_0,\phi,\psi}(x_s,s\mid o,c)\\
  {}-(a-x_0)\bigr\|_2^2\Bigr].
\end{split}
\label{eq:fm}
\end{equation}
При генерации интегрирование векторного поля от $s=0$ к $s=1$ преобразует шум в блок управляющих действий.

Для замороженного линейного слоя с весовой матрицей $W\in\mathbb{R}^{m\times n}$ исходное преобразование дополняется низкоранговой поправкой~\cite{lora}:
\begin{equation}
  h=Wx+\frac{\alpha}{r}W_{\mathrm{up}}W_{\mathrm{down}}x.
\label{eq:lora}
\end{equation}
Здесь $W_{\mathrm{down}}\in\mathbb{R}^{r\times n}$ проецирует вход в подпространство малой размерности $r$, а $W_{\mathrm{up}}\in\mathbb{R}^{m\times r}$ возвращает его в размерность выхода; $r\ll\min(m,n)$. Матрица $W$ остается неизменной. Множитель $\alpha/r$ задает масштаб поправки.

Полная настройка слоя потребовала бы переобучения $mn$ элементов, тогда как обычный LoRA-адаптер содержит $(m+n)r$ параметров. Отношение
\begin{equation}
  \frac{(m+n)r}{mn}
\end{equation}
мало при достаточно низком ранге. В гиперсетевой схеме эти матрицы порождаются, поэтому число оптимизируемых параметров всей системы определяется размером гиперсети и эксперта действий, а не только размером адаптера.

Гиперсеть устроена как композиция энкодера контекста и выходных голов. Сначала контекст $c$ отображается в скрытый эмбеддинг $e=f_\phi(c)$, затем головы $g_\phi$ порождают параметры адаптера. С учетом индекса слоя $l$ получаем
\begin{equation}
  \bigl(W_{\mathrm{down}}^{(l)},W_{\mathrm{up}}^{(l)}\bigr)
  =g_\phi\bigl(f_\phi(c),l\bigr).
\label{eq:hyper}
\end{equation}
При совпадающих размерностях адаптируемых матриц одни и те же выходные головы переиспользуются для разных слоев; для разных типов проекций необходимы совместимые по размерности головы. Такое разделение сокращает стоимость по сравнению с независимым полным генератором на каждый слой~\cite{hyperlora}.

Энкодер контекста может быть реализован многослойным перцептроном либо трансформером. В трансформерном варианте инструкция и представления индексов адаптируемых слоев подаются как токены единой последовательности, а внимание согласует контекст задачи с представлением конкретного слоя при порождении его весов.

В настоящей постановке адаптер порождается для каждого примера батча. На каждом прямом проходе по инструкции генерируются собственные матрицы $W_{\mathrm{down}}$ и $W_{\mathrm{up}}$, которые используются в MLP-слоях языковой части VLM. Градиент функции потерь~\eqref{eq:fm} проходит через замороженные слои и обновляет $\phi$ и $\psi$, не изменяя $\theta_0$.

Такое сквозное обучение проводится на фиксированных демонстрациях. Порождение весов внутри прямого прохода не предполагает обучение с подкреплением или сбор новых взаимодействий. Замораживание параметров VLM также не означает отключения вычисления градиента по входящим в него адаптерам.

Способность гиперсети порождать осмысленные адаптеры для невиданных условий зависит от богатства контекста и риска переобучения. При бедном контексте выразительный генератор может запоминать соответствие «обучающая инструкция -- комплект весов». Для снижения этого риска предусматривается dropout контекстного эмбеддинга перед выходными головами~\cite{hyperlora}.

В расширенной конфигурации текстовый контекст предполагается дополнять визуальными признаками текущего наблюдения, извлекаемыми замороженным энкодером DINOv2~\cite{dinov2}. Тогда $c=(c_{\mathrm{text}},c_{\mathrm{visual}})$, и адаптер зависит не только от формулировки задачи, но и от наблюдаемой сцены.

Наличие вариаций в обучающем распределении не гарантирует перенос, а отсутствие прямых примеров определенного сдвига не доказывает его невозможность. Поэтому поведение по осям object, position, semantic и task рассматривается как предмет эмпирической проверки. Качество планируется представлять отдельно по каждой оси, без замены этих результатов единым средним показателем.

\Needspace{5\baselineskip}
\begin{thebibliography}{99}

\bibitem{firoozi2023}
Firoozi R., Tucker J., Tian S. et al. Foundation Models in Robotics: Applications, Challenges, and the Future. 2023. arXiv:2312.07843. \url{https://arxiv.org/abs/2312.07843}.

\bibitem{rt1}
Brohan A., Brown N., Carbajal J. et al. RT-1: Robotics Transformer for Real-World Control at Scale. 2022. arXiv:2212.06817. \url{https://arxiv.org/abs/2212.06817}.

\bibitem{rt2}
Brohan A., Brown N., Carbajal J. et al. RT-2: Vision-Language-Action Models Transfer Web Knowledge to Robotic Control. 2023. arXiv:2307.15818. \url{https://arxiv.org/abs/2307.15818}.

\bibitem{oxe}
Open X-Embodiment Collaboration. Open X-Embodiment: Robotic Learning Datasets and RT-X Models. 2023. arXiv:2310.08864. \url{https://arxiv.org/abs/2310.08864}.

\bibitem{octo}
Octo Model Team et al. Octo: An Open-Source Generalist Robot Policy. 2024. arXiv:2405.12213. \url{https://arxiv.org/abs/2405.12213}.

\bibitem{openvla}
Kim M. J., Pertsch K., Karamcheti S. et al. OpenVLA: An Open-Source Vision-Language-Action Model. 2024. arXiv:2406.09246. \url{https://arxiv.org/abs/2406.09246}.

\bibitem{smolvla}
Shukor M., Aubakirova D., Capuano F. et al. SmolVLA: A Vision-Language-Action Model for Affordable and Efficient Robotics. 2025. arXiv:2506.01844. \url{https://arxiv.org/abs/2506.01844}.

\bibitem{liberopro}
Zhou X., Xu Y., Tie G. et al. LIBERO-PRO: Towards Robust and Fair Evaluation of Vision-Language-Action Models Beyond Memorization. 2025. arXiv:2510.03827, v2, 2026. \url{https://arxiv.org/abs/2510.03827v2}.

\bibitem{libero}
Liu B., Zhu Y., Gao C. et al. LIBERO: Benchmarking Knowledge Transfer for Lifelong Robot Learning. \textit{Advances in Neural Information Processing Systems}. 2023. Vol. 36. P. 44776--44791. \url{https://arxiv.org/abs/2306.03310}.

\bibitem{lora}
Hu E. J., Shen Y., Wallis P. et al. LoRA: Low-Rank Adaptation of Large Language Models. \textit{ICLR}. 2022. \url{https://arxiv.org/abs/2106.09685}.

\bibitem{hyperlora}
Sharma S., Xiong Z., Li K., Vuorio R., Whiteson S. Efficient Domain Adaptation of Robotic Foundation Models via Hypernetwork-Generated LoRA. \textit{NeurIPS 2024 Workshop}. 2024. \url{https://openreview.net/forum?id=SEvvWs3CAL}.

\bibitem{simpler}
Li X., Hsu K., Gu J. et al. Evaluating Real-World Robot Manipulation Policies in Simulation. 2024. arXiv:2405.05941. \url{https://arxiv.org/abs/2405.05941}.

\bibitem{flowmatching}
Lipman Y., Chen R. T. Q., Ben-Hamu H. et al. Flow Matching for Generative Modeling. \textit{ICLR}. 2023. \url{https://arxiv.org/abs/2210.02747}.

\bibitem{dinov2}
Oquab M., Darcet T., Moutakanni T. et al. DINOv2: Learning Robust Visual Features without Supervision. 2023. arXiv:2304.07193. \url{https://arxiv.org/abs/2304.07193}.

\bibitem{liberodata}
Physical Intelligence. LIBERO: LeRobot dataset (physical-intelligence/libero). Карточка данных. \url{https://huggingface.co/datasets/physical-intelligence/libero}. Дата обращения: 07.11.2026.

\end{thebibliography}
\end{document}
